% from ega4-1-fr.tex
\oldpage[IV-1]{227}
(où $\pi$ est le morphisme structural) est commutatif et identifie
$X\times_Y X$ au produit des $(Y\times_Z Y)$-préschémas $Y$ et
$X\times_Z X$ (I, 5.3.5). Si $g$ est quasi-séparé, $\Delta_g$ est
quasi-compact, donc il en est de même de $i$ (1.1.2, (iii))~; si de plus $f$
est quasi-séparé, $\Delta_f$ est quasi-compact, donc il en est de même de
$i\circ\Delta_f$ (1.1.2, (ii)) qui est égal à $\Delta_{g\circ f}$.
\end{proof}

\begin{corollary}[1.2.3]
\label{IV.1.2.3-fr}
\begin{enumerate}
  \item[(i)] Soit $X$ un préschéma quasi-séparé. Alors tout morphisme $f:X\to Y$
  est quasi-séparé.
  \item[(ii)] Soit $Y$ un préschéma quasi-séparé. Pour qu'un morphisme $f:X\to Y$
  soit quasi-séparé, il faut et il suffit que le préschéma $X$ soit
  quasi-séparé.
\end{enumerate}
\end{corollary}

Cela résulte aussitôt de (1.2.2, (ii) et (v)) appliqué à $X$, $Y$ et
$\operatorname{Spec}(\mathbf Z)$.

\begin{proposition}[1.2.4]
\label{IV.1.2.4-fr}
Soient $f:X\to Y$ un morphisme, $g:Y\to Z$ un morphisme quasi-séparé. Si
$g\circ f$ est quasi-compact, il en est de même de $f$.
\end{proposition}

\begin{proof}
On sait que $f$ est le morphisme composé
$X\xrightarrow{\Gamma_f}X\times_ZY\xrightarrow{\mathrm{pr}_2}Y$
(I, 5.3.13)~; d'autre part, $\mathrm{pr}_2$ s'identifie à
$(g\circ f)\times_Z1_Y$ (I, 3.3.4) et si $g\circ f$ est quasi-compact, il en
est de même de $\mathrm{pr}_2$ (1.1.2, (iv)). Enfin, on a le diagramme
commutatif
\[
\xymatrix{
  X \ar[r]^-{\Gamma_f} \ar[d]_-{f} & X\times_ZY
    \ar[d]^-{f\times 1_Y} \\
  Y \ar[r]_-{\Delta_g} & Y\times_ZY
}
\tag{1.2.4.1}
\label{IV.1.2.4.1-fr}
\]
qui identifie $X$ au produit des $(Y\times_ZY)$-préschémas $Y$ et
$X\times_ZY$ (I, 5.3.7). Comme par hypothèse $\Delta_g$ est un morphisme
quasi-compact, il en est de même de $\Gamma_f$ (1.1.2, (iii))~; la conclusion
résulte donc de (1.1.2, (ii)).
\end{proof}

\begin{proposition}[1.2.5]
\label{IV.1.2.5-fr}
Soient $f:X\to Y$, $g:Y'\to Y$ deux morphismes,
$X'=X\times_Y Y'$, $f'=f_{(Y')}:X'\to Y'$. Si $g$ est quasi-compact et
surjectif et $f'$ quasi-séparé, alors $f$ est quasi-séparé.
\end{proposition}

\begin{proof}
On a en effet $(X'\times_{Y'}X')=(X\times_YX)_{(Y')}$ et
$\Delta_{f'}=(\Delta_f)_{(Y')}$ (I, 5.3.4)~; comme la projection
$X'\times_{Y'}X'\to X\times_YX$ est quasi-compacte (1.1.2, (iii)) et
surjective (I, 3.5.2), on peut, en vertu de (I, 3.3.11), appliquer (1.1.4),
qui montre que puisque $\Delta_{f'}$ est un morphisme quasi-compact, il en est
de même de $\Delta_f$.
\end{proof}

\begin{proposition}[1.2.6]
\label{IV.1.2.6-fr}
Soient $f:X\to Y$ un morphisme, $(U_\alpha)$ un recouvrement de $Y$ par des
ouverts tels que les préschémas induits $U_\alpha$ soient quasi-séparés. Pour
que $f$ soit quasi-séparé, il faut et il suffit que chacun des préschémas
induits par $X$ sur les ouverts $f^{-1}(U_\alpha)$ soit quasi-séparé.
\end{proposition}

\begin{proof}
L'image réciproque dans $X\times_YX$ de $U_\alpha$ est
$X_\alpha\times_{U_\alpha}X_\alpha$, en posant
$X_\alpha=f^{-1}(U_\alpha)$, et la restriction
$X_\alpha\to X_\alpha\times_{U_\alpha}X_\alpha$ de $\Delta_f$ n'est autre
que $\Delta_{f_\alpha}$, en désignant par $f_\alpha$ la restriction
$X_\alpha\to U_\alpha$ de $f$~; en vertu de la définition (1.2.1) et du
caractère local de la notion de morphisme quasi-compact (1.1.1), pour que $f$
soit quasi-séparé, il faut et il suffit que chacun des morphismes $f_\alpha$
le soit. Mais comme par hypothèse le morphisme
$U_\alpha\to\operatorname{Spec}(\mathbf Z)$ est quasi-séparé, il revient au
même de dire que $f_\alpha$ est quasi-séparé ou que le composé
$X_\alpha\xrightarrow{f_\alpha}U_\alpha\to\operatorname{Spec}(\mathbf Z)$
l'est (1.2.2, (ii) et (v)), d'où la proposition.
\end{proof}

